\begin{table}
\begin{tabular}{ | m{5em} | m{17em} | } 
  \hline
  Frontend & 2-way superscalar, 256 entries BTB, 8 entries RAS, tournament predictor \\
  \hline
  ROB &  64 entries, 2-way insert/commit \\
  \hline
  Reservation Stations & 2 stations for ALU (16 entries), 1 for FPU (16 entry), 1 for memory (16 entries) \\
  \hline
  L1 (I/D) &  32KB, 8 ways, 8 outstanding requests \\
  \hline
  Ld-St Unit & 24 LdQ entries, 14 StQ entries, 4 SB entries\\
  \hline
  LLC (L2) & 1MB, 10 cycles latency, 16 ways, max 16 outstanding req. \\
  \hline
  Memory &  120 cycles latency (24 outstanding reqs) \\
  \hline
\end{tabular}
    \caption{Citadel's microarchitectural configuration.}
    \label{tab:mi6param}
\end{table}
\begin{table}
    \centering
    \begin{tabular}{|p{3cm}|r|r|}
        \hline
         Component & LOC & Size\\
         \hline\hline
         \rowcolor{LightRedbis}
         Processor\cite{bourgeat2019mi6} (+mod) & $\sim60$ KLOC (+ 350) & N/A \\
         \hline
         \rowcolor{myblue}
         Bootloader & 5236 & 55KB\\
         \hline
         \rowcolor{myblue}
         Security Monitor & 9287 & 97KB\\
         \hline
         \rowcolor{LightRedbis}
         Mini-SM & 2301 & 5.6KB\\
         \hline
         \rowcolor{LightRed}
         Linux & $\sim20$M & 157MB\\
         \hline
         \rowcolor{LightRed}
         SM Kernel Mod. & 507 & 505KB\\
         \hline
         \hline
         \rowcolor{myblue}
         TCB Software & 14523 & 152KB \\
         \hline
    \end{tabular}
    \caption{Software Component Size and TCB Breakdown. Mini-SM is a subset of the SM code. 
    Lines in red correspond to elements excluded from the TCB.}
    \label{tab:tcbloc}
    \smallskip
    \centering
    \small
    \begin{tabular}{|p{1.2cm}|r|r|}
        \hline
         Modules & LUTs & FFs\\
         \hline\hline
         Top & 410K (\textcolor{blue}{+3.0\%} | \textcolor{orange}{+1.8\%}) & 251K (\textcolor{blue}{+4.4\%} | \textcolor{orange}{+2.6\%})\\
         \hline
         Core & 169K (\textcolor{blue}{+1.6\%} | \textcolor{orange}{+0.3\%}) & 91K (\textcolor{blue}{+3.4\%} | \textcolor{orange}{+1.0\%})\\
         \hline
    \end{tabular}
    \caption{Hardware overhead for the top module and the core, lookup tables (LUTs) and flip flops (FFs). First comparison is to Riscy-OOO while second is to MI6 with the parameters of Table \ref{tab:mi6param}.}
    \label{tab:hwoverhead}
\end{table}

\section{Evaluation}
\label{sec:eval}

\noindent\textbf{Experimental Setup}
Citadel's processor is a modified out-of-order, superscalar, 2-core based on Riscy-OOO~\cite{zhang2018composable} and MI6~\cite{bourgeat2019mi6}.
The parameters used for the cores and the memory subsystem are given in Figure~\ref{tab:mi6param}, and kept small to keep the synthesis and simulation time small.
Citadel runs on an AWS EC2 F1 instance, equipped with a Virtex UltraScale+ FPGA. 
On this FPGA board, Citadel can be clocked at $30$ MHz, which is similar to the baseline design.
We use the RISC-V tool chain with GCC 12.2.0.
We compile all software with the -O3 flag. 
We run Linux Kernel 6.2.0.
All numbers are collected by running the benchmarks on the FPGA using performance counters on an average of ten runs. 
\\

\noindent\textbf{Implementation Data, TCB and Surface Overhead}
Number of lines of code (LOC) and binary size of Citadel's components are presented in Table \ref{tab:tcbloc}.
Due to their limited functions, the SM and the Bootloader can be kept small and the total software Trusted Code Base (TCB) is less than 15K LOC.
This helps minimize potential bugs and reduce the software attack surface.
It can also help make future formal verification efforts tractable.
The final TCB size needs to include the enclave's binary, which depends on the application.
Table \ref{tab:hwoverhead} shows the area overhead for Citadel compared to Riscy-OOO\cite{zhang2018composable} and  MI6\cite{bourgeat2019mi6}.
Most of the overhead comes from enforcing microarchitectural isolation (i.e. MI6) or from the modifications in the LLC with, compared to Mi6, the addition of an arbiter, a fine-grained partitioning mechanism and the L1 bypass mechanism.
The overhead for the core compared to MI6 is minimal (+0.3\% LUTs and +1.0\% FFs), which gives us an upper bound for Safe mode and Burst mode.
This is significantly lower than other mechanisms for secure speculation (see Section \ref{sec:relatedwork}).


\begin{table}[ht]
    \centering
    \begin{tabular}{|c|c|c|}
    \cline{2-3}
    \multicolumn{1}{c|}{} & \textbf{Memcpy} & \textbf{Random} \\ \hline
    \textbf{Baseline} & 343928  &  273164 \\ \hline
    \textbf{Safe} & 595362 (+73\%) & 1447543 (+430\%) \\ \hline
    \rowcolor{mygreen!35}
    \textbf{Safe+Burst} & 343291 (+0\%) & 270394 (-1\%) \\ \hline
    \textcolor{red}{\textbf{Safe+Burst+L1}} & 344163 (+0\%) & 277969 (+2\%) \\ \hline
    \end{tabular}
    \caption{Microbenchmarks for \texttt{memcpy} and random accesses to shared memory. \texttt{memcpy} copies 128MB from shared memory to private memory. Baseline is \texttt{memcpy} from an to private memory with speculation enabled. \emph{Random} accesses 10240 elements in a 128MB memory region.
    We experiment with \emph{insecurely} enabling the L1 for shared memory accesses.}
    \label{tab:resultmemcpy}
\newcolumntype{C}[1]{>{\centering\arraybackslash}m{#1}}
     \begin{NiceTabular}{|C{2.7cm}|C{1.2cm}|C{1.2cm}|C{1.4cm}|}[hvlines, corners=NW]
        \CodeBefore
        \rowcolor{gray!35}{2}
        \rowcolor{mygreen!35}{4}
        \Body
            \multicolumn{1}{c|}{} & \textbf{ECC} & \textbf{ML} & \textbf{ML EXT}\\
            \Block{}{\textcolor{red}{\textbf{{Baseline}}}} & 98599278 & 406765555 & 406765555 \\
            \Block{}{\textbf{Safe Only}} & 102433204 (+4\%) & 419853901 (+3.2\%) & 1119486865 (+175\%)  \\
            \Block{}{\textbf{Citadel\\(Safe+Burst)}} & 100411982 (+1.8\%) & 421332976 (+3.5\%) & 426527197 (+4.8\%) \\
    \end{NiceTabular}
    \caption{Performance overheads for applications that access shared memory. Our elliptic curve cryptography (ECC) benchmark signs a total of 1MB of data split in 3000 packets. Our ML benchmark evaluate private inference over 1024 encrypted requests. ML EXT is the same but the public model weights are in shared memory.
    Requests are sent by an untrusted app using message passing through shared memory. 
    Insecure Baseline is the same application directly linked into the untrusted app.
    }
    \label{tab:resultcrypto}
\end{table}

\subsection{Microbenchmarks}
\label{sec:evalsharedmemory}

We evaluate two \emph{pathological} microbenchmarks with distinct but shared memory access patterns: \emph{Memcpy} (sequential access, copying 128KB from shared to enclave private memory) and \emph{Random} (10,240 random accesses within a 128KB region).
Results can be found in Table \ref{tab:resultmemcpy}.
We compare these with a baseline of private memory accesses.
For all tests, we disable LLC partitioning and use 4KB page mapping.
Safe mode alone significantly increases overhead (+73\% for Memcpy, +430\% for Random) due to disabling the pipelining of shared memory accesses. 
Activating Burst mode for both access types (after verifying security with our code analysis tool) dramatically reduces this overhead, approaching baseline performance.
Note that verification is a one-time task and that Burst-enable \emph{memcopy} can now be used by other code \emph{without further program analysis}.
We also experiment with (insecurely) disabling the L1 bypass for shared memory accesses, which, for these benchmarks, has negligible performance impact. 
Indeed, for \texttt{memcpy}, adjacent pipelined loads will always be resolved on the next cycle (in the L1 or the LLC) and similarly than for Random, data is not re-accessed.
For Random, these results show that polluting the L1-D with random memory accesses slightly degrades performance.

\subsection{Evaluating the End-to-End Platform}

We evaluate the platform usability by porting existing applications, showing how Burst mode can be easily leveraged to gain performance at the cost of minimal code annotation and analysis.
Our shared memory mechanisms only incur performance overheads when the enclave program accesses shared memory, making evaluation mostly relevant for programs with private and shared memory (e.i. not SPEC).
Note that for non-enclave programs, the only overhead is due to the LLC arbiter and is minimal on our two-cores prototype. We measure $\approx$0.1\% overhead compared to Riscy-OOO~\cite{zhang2018composable} on our different benchmarks when run as insecure baselines. 
\\

\noindent\textbf{Securing a Crypto Library (ECC)}
\label{sec:evalcrypto}
We write a wrapper around an ED25519 RISC-V library \cite{ed25519github}.
We add a queue to shared memory, so untrusted applications can send requests to the library to generate keys (that will stay in enclave memory), or to sign messages using previously generated keys.
In total, we need to implement 334 LOC of C, which is to be compared to the 4,514 LOC of the crypto library (7\% overhead).
With the enclave's code added, the total TCB is now 19,231 LOC.
The final crypto-enclave binary is 71KB.
As shown in Table \ref{tab:resultcrypto}, we see minimal overhead mainly due to message passing and copying inputs inside the enclave.
We instrument the SHA256 macro that performs the input copy to leverage Burst mode.
We successfully verify that the code snippet is secure using our analysis tool we reduce the overhead to 1.8\%.
For completeness, and because the ECC benchmark is implemented using constant-time programming \cite{reparaz2017dude, almeida2016verifying}, we also evaluate the performance overhead of disabling speculation completely in order to defend against "execution time" attacks (see Section \ref{sec:execution_time_SC}). We observe a +882\% overhead, orders of magnitude above our other mechanisms, which clearly highlights the trade-off Citadel offers between security and performance.
\\

\noindent\textbf{Private Inference (ML and ML EXT)}
\label{sec:evalprivateinference}
We convert a handwritten digit recognition model (trained on MNIST) to standard C code \cite{onnx2cgithub} and run private inference inside our enclave.
A user uses our remote attestation mechanism to verify the enclave identity and establish a secure communication channel.
When receiving requests through shared memory, the enclave decrypts the ciphertext, performs the inference task on the private input and sends back an encrypted output.
We also experiment with placing the model weights in shared memory (ML EXT).
This scenario is especially relevant as some models might be too big to fit in DRAM and require to leverage demand paging. 
Note that this type of access patterns would be hard to reproduce using message-passing primitives (e.g. DMA accesses).
In this scenario, the integrity of the model weights is not protected.
In total, we need to implement 227 lines of straightforward C, which is to be compared to the 1536 LOC of the neural network, the 410 LOC of the AES library and the 4386 LOC of ED25519 we also include
for encryption and key exchange.
The TCB is now 21K LOC and the final enclave binary measures 54KB.
Setting up the enclave, represents a (one-time) cost similar to performing 235 inferences.
For a given request, runtime is split 55\% and 45\% between AES and inference.
Table \ref{tab:resultcrypto} show that overheads are reasonable but can get up to +175\% when placing the model in shared memory.
We use our Burst-instrumented \texttt{memcpy} function but see no performance improvement as only small messages are copies.
Nevertheless, on ML EXT, we instrument the inference code to leverage Burst (+2 LOC).
Our code analysis tool revealed a Spectre-like gadget as the pointer to the tensor is first dereferenced before the tensor values are accessed.
We made sure to deference the first pointer before turning on Burst mode (+1 LOC), enough to pass program analysis.
We observe performance overhead drops down to 4.8\%. 
These results highlight the versatility of shared memory and the tradeoff between performance and requiring program analysis.
Finally, insecurely disabling L1-bypass for shared memory on ML EXT only show modest performance gains (+4.3\% down from 4.8\% ).
\\ 

\noindent\textbf{MicroPython Runtime}
\label{sec:evalmicropython}
We port MicroPython \cite{microphython} inside an enclave
which serves as a small benchmark to demonstrate the expressivity and flexibility of the platform.
We redirect the console to shared memory and a hash of the standard output is extended on each character print.
The hash of the console transcript is signed on the enclave exit using ECC and the enclave's own private keys (obtained through the local attestation mechanism).
In total, we modify 128 LOC of MicroPython.
With the enclave's code added, the TCB is now 65K LOC, most of it (70\%) being MicroPython.
We successfully run our attested runtime on a simple primality test written in Python.

\subsection{Evaluating LLC Dynamic Partitioning}
\label{sec:evalllcpartitioning}

We evaluate dynamically partitioning the LLC on our cryptographic enclave benchmark.
We compare 1)
an insecure baseline where the LLC is not partitioned, 2) a secure configuration where the LLC is uniformly statically partitioned (MI6 scheme) and 3) a non-uniform partitioning set up dynamically using our scheme.
For the non-uniform configuration, we split the 1MB of LLC as follows: we give 256KB to the enclave, 16KB to the SM, 128KB to the OS, 128KB to the region used for shared memory, 128KB to the region storing page tables, and 1KB each to every other region.
We observe a 67.8\% overhead for the uniform partitioning and 0.08\% overhead for the non-uniform partitioning.
This is especially important as untrusted applications also benefit from it.
Modifying LLC partitioning (infrequent) takes on average 500K cycles.
Among these, 98\% of the cycles are spent flushing the LLC.
Resizing fewer regions is faster: we measure an average of 93 cycles to flush a single LLC set containing 16 ways.